\begin{lemma}{Integrierbarkeit und Verteilung von Linearkombinationen}
             {IntegrierbarkeitUndVerteilungVonLinearkombinationen}
Seien $t \in \R^m$, $v \in \R^N$, $c \in \R$.\\
Dann gilt:
\begin{enumerate}[(1)]
\item\label{IntegrierbarkeitUndVerteilungVonLinearkombinationen1}
$c\bprod t \delta_m + \bprod v\gamma_N, \, 
\sprod \cdot v_N \circ ev(t, \cdot) \circ g \in L_2(\Omega,\mfa,\Prim,\R)$.
\item\label{IntegrierbarkeitUndVerteilungVonLinearkombinationen2}
$c\bprod t \delta_m + \bprod v\gamma_N = 0 \iff (ct = 0 \wedge v = 0)$.
\item\label{IntegrierbarkeitUndVerteilungVonLinearkombinationen3}
$\sprod \cdot v_N \circ ev(t, \cdot) \circ g = 0 \iff (t = 0 \vee v= 0)$.
\item\label{IntegrierbarkeitUndVerteilungVonLinearkombinationen4}
Ist $ct \neq 0$ oder $v \neq 0$, so gilt:
\[c\bprod t \delta_m + \bprod v\gamma_N \quad{ist}\quad
N\left(0, c^2\norm t_{2,m}^2 +\norm v_{2, N}^2\right)\text{--verteilt.}\]
\item\label{IntegrierbarkeitUndVerteilungVonLinearkombinationen5}
Ist $t \neq 0\neq v $, so gilt:
$\sprod \cdot v_N \circ ev(t, \cdot) \circ g$ ist
$N\left(0, \norm t _{2,m}^2\norm v _{2, N}^2\right)$--verteilt.
\end{enumerate}
\end{lemma}
\begin{proof}
(1) $c\bprod t \delta_m + \bprod v\gamma_N$ ist als Linearkombination
quadratintegrierbarer Funktionen selbst quadratintegrierbar.\\
Ferner gilt f�r alle $\omega \in \Omega$:
\[\bigg(\sprod \cdot v_N \circ ev(t, \cdot) \circ g\bigg)(\omega) =
\sprod {g(\omega)t}v_N = \sum_{i=1}^N\sum_{j=1}^mt_jv_i(g_{ij}(\omega))\text.\]
Damit ist\[
\sprod\cdot v_N\circ ev(t,\cdot)\circ g=\sum_{i=1}^N\sum_{j=1}^mt_jv_ig_{ij}
\text,\]
also ebenfalls als Linearkombination quadratintegrierbarer Abbildungen
quadratintegrierbar.\\
(2) Da $\gamma, \delta$ vollst�ndig unabh�ngig und zentriert sind, gilt
\[\forall\,i\in\haken m\,\forall\,j\in\haken N:\E(\delta_i\gamma_j)=0\text.\]
Da $\E$ das Skalarprodukt auf $L_2(\Omega, \mfa, \Prim, \R)$ erzeugt, sind also
die Komponentenfunktionen von $\gamma$ und $\delta$ paarweise orthogonal und
damit linear unabh�ngig. Damit gilt:
\begin{align*}
c\bprod t \delta _m + \bprod v \gamma_N = 0 
& \iff \left(\sum_{i=1}^mct_i\delta_i \right) 
     + \left(\sum_{j=1}^Nv_j\gamma_j \right) = 0\\
& \iff(\forall\,i\in \haken m: ct_i = 0) \wedge (\forall\,j\in\haken N:v_j=0)\\
& \iff ct = 0 \wedge v = 0\text.
\end{align*}
(3) Weil $\widetilde g$ unabh�ngig und zentriert ist, sind die
Komponentenfunktionen von $\widetilde g$ und damit auch jene von $g$ paarweise
orthogonal, also linear unabh�ngig. Somit folgt wie im Beweis von (1):
\begin{align*}
 \sprod \cdot v_N \circ ev(t, \cdot) \circ g 
 \overset{\substack{\text{Bew.}\\ \text{von }(1)}}=
  \sum_{i=1}^N\sum_{j=1}^mt_jv_ig_{ij} = 0
&\iff \forall\, i \in \haken N \forall\, j \in \haken m: t_jv_i = 0\\
& \iff  t= 0 \vee v = 0\text.
\end{align*}
(4)/(5) Sind $X_1,X_2$ stochastisch unabh�ngige Zufallsgr��en, $a_1,a_2\in \R$
und $\sigma_1,\sigma_2\in\R_{>0}$ derart, dass $X_i$
$N(a_i,\sigma_i^2)$--verteilt sind f�r alle $i \in \haken 2$, so ist nach \cite{Irle} (dort 10.15) $X_1 + X_2$
$N(a_1+a_2, \sigma_1^2 + \sigma_2^2)$--verteilt.\\
Seien $c\in \R\setminus\meng{0}$ und $t \in \R$. Dann gilt falls $c > 0$:
\begin{align*}
\Prim(cX_1 \le t) &= \Prim(X_1 \le \nicefrac t c) 
= \frac 1{\sqrt{2\pi\sigma_1^2}}\int_{-\infty}^{\nicefrac t c}
\exp\left(- \frac 1 2 \left(\frac{s-a_1}{\sigma_1}\right)^2 \right)ds\\
&\overset{\substack{\text{Subst.}\\s \mapsto \frac s c}}=
\frac 1{\sqrt{2\pi\sigma_1^2}}\int_{-\infty}^t 
\exp\left(-\frac 1 2 \left(\frac{\frac s c - a_1}{\sigma_1} \right)^2
 \right)\frac 1 c ds\\
 &=\frac 1 {\sqrt{2 \pi (c\sigma_1)^2}}\int_{-\infty}^t 
 \exp\left(-\frac 12\left(\frac{s-(ca_1)}{(c\sigma_1)}\right)^2 \right)ds
\end{align*}
und falls $c < 0$ wegen der Stetigkeit von $N(a_1, \sigma_1^2)$:
\begin{align*}
&\,\Prim(cX_1 \le t) = \Prim(X_1 \ge \nicefrac tc) 
= 1 - \Prim(X_1 \le \nicefrac t c)\\
=&\,
\frac 1 {\sqrt{2\pi\sigma_1^2}}\int_{\frac t c}^\infty
\exp\left(- \frac 1 2 \left(\frac{s-a_1}{\sigma_1}\right)^2 \right)ds\\
\overset{\substack{\text{Subst.}\\s \mapsto \frac s c}}=&\,
\frac 1 {\sqrt{2\pi\sigma_1^2}}\int_{ t }^{-\infty}
\exp\left(-\frac 12\left(\frac{\frac sc-a_1}{\sigma_1}\right)^2\right)
\frac 1 c ds\\
=&\, \frac 1{\sqrt{2\pi\sigma_1^2}}\frac 1{\abs c}
\int_{-\infty}^t \exp\left(-\frac 12\left(\frac{ s-(ca_1)}{c\sigma_1}\right)^2\right)ds\text.
\end{align*}
Also ist $cX_1$ auch $N(ca_1, (c\sigma_1)^2)$--verteilt.\\
Damit folgen die Aussagen (4) und (5) durch einfache Induktionen.
\end{proof}